% GENERATED FILE. Edit canonical lessons, then run npm run generate:books.
% canonical-source-hash: fnv1a64:31898bd4820a6171
% canonical-lessons: IT-C03-io, IT-C03-mi-chiamo, IT-C03-come-ti-chiami, IT-C03-piacere, IT-C03-practice

\chapter{Introducing Yourself}
\label{ch:it-introductions}
\hlchaptermodality{3}

\begin{chapteropening}
\textbf{By the end of this chapter:} \emph{I can give my name in Italian, ask for someone else's, and say that it is a pleasure.}

A whole first meeting: ciao, mi chiamo …, Come ti chiami? in the informal and Come si chiama? in the formal, and piacere to close it.
\end{chapteropening}

\section[io]{io — \textquotedblleft{}I,\textquotedblright{} the pronoun Italian loves to leave out}

\label{lesson:IT-C03-io}



\begin{quote}
Before you introduce yourself, meet \textbf{io} — \textquotedblleft{}I.\textquotedblright{} It's a small word with a famous English cousin, and a surprising habit: Italian usually \textbf{drops it.}
\end{quote}

\begin{sounds}
\begin{itemize}
\raggedright
  \item \texttt{io-glide} — \textbf{io} = \emph{EE-oh}, two quick vowels gliding together, stress the first: \emph{EE-oh}.
\end{itemize}
\end{sounds}

\begin{cousinweb}
\textbf{io} = \textquotedblleft{}I,\textquotedblright{} worn down from Latin \textbf{ego} — which English borrowed \emph{whole}:

\begin{itemize}
\raggedright
  \item \textbf{ego} — literally Latin \textquotedblleft{}I\textquotedblright{} (your sense of self).
  \item \textbf{egoist}, \textbf{egotist}, \textbf{egocentric} — all built on that \textquotedblleft{}I.\textquotedblright{}
  \item Its Romance siblings: Spanish \textbf{yo}, Portuguese \textbf{eu}, French \textbf{je} — every one a child of \emph{ego}.
\end{itemize}
\end{cousinweb}

\begin{grammarlens}[title={Grammar Lens: the dropped subject (pro-drop)}]
Here's the Italian habit worth learning now: \textbf{the verb ending already tells you who's speaking, so the pronoun is usually left out.} \emph{Sto bene} (\textquotedblleft{}I'm well\textquotedblright{}) needs no \emph{io} — the \emph{-o} of \emph{sto} is \textquotedblleft{}I.\textquotedblright{} You add \textbf{io} only for \textbf{emphasis} or \textbf{contrast}:

\begin{quote}
\emph{Io} sto bene — e tu? — \textquotedblleft{}\emph{I'm} fine — and you?\textquotedblright{}
\end{quote}

Spanish (\emph{yo}) and Portuguese (\emph{eu}) do exactly the same. English can't: it must say \textquotedblleft{}I\textquotedblright{} every time. So learning Italian means learning to \emph{trust the verb ending} and let the pronoun go.
\end{grammarlens}

\subsection*{Your turn}
Say these aloud:

\begin{itemize}
\raggedright
  \item \textquotedblleft{}io\textquotedblright{} — \emph{EE-oh}
  \item \textquotedblleft{}io\textquotedblright{} then English \textquotedblleft{}ego, egoist\textquotedblright{} — the same Latin \textquotedblleft{}I\textquotedblright{}
  \item \textquotedblleft{}Sto bene\textquotedblright{} with no \emph{io} — the ending carries the \textquotedblleft{}I\textquotedblright{}
\end{itemize}

\begin{tcolorbox}[breakable,colback=teal!4,colframe=teal!35!black,title={Before you move on}]
What Latin word is \emph{io} from, and what English words share it? (\emph{Ego} — ego, egoist.) Why does Italian usually drop \emph{io}? (The verb ending already says who.) When do you keep it? (For emphasis/contrast.) Next: use it to introduce yourself — \textbf{mi chiamo}.
\end{tcolorbox}
\section[mi chiamo]{mi chiamo — \textquotedblleft{}I call myself\textquotedblright{}}

\label{lesson:IT-C03-mi-chiamo}



\begin{quote}
Like Spanish and French, Italian doesn't say \textquotedblleft{}my name is.\textquotedblright{} It says \textbf{\textquotedblleft{}I call myself\textquotedblright{}} — \textbf{mi chiamo}. Same reflexive trick, same \emph{clāmāre} root as Spanish \emph{me llamo}.
\end{quote}

\begin{sounds}
\begin{itemize}
\raggedright
  \item \texttt{ch-hard-k} — Italian \textbf{ch} is a \textbf{hard k} (the opposite of English!): so \textbf{chiamo} = \emph{KYAH-moh}, not \textquotedblleft{}chee-.\textquotedblright{} (Italian softens \emph{c} to \emph{ch}-sound only before \emph{e/i}; the \emph{h} in \emph{chi} is there precisely to keep it hard.)
  \item \textbf{mi chiamo} = \emph{mee KYAH-moh}.
\end{itemize}
\end{sounds}

\begin{cousinweb}
\textbf{mi chiamo} = \textbf{mi} (\textquotedblleft{}myself\textquotedblright{}) + \textbf{chiamo} (\textquotedblleft{}I call\textquotedblright{}) $\to$ \textbf{\textquotedblleft{}I call myself.\textquotedblright{}} Then the name: \emph{Mi chiamo Marco} — \textquotedblleft{}I call myself Marco.\textquotedblright{}

The verb is \textbf{chiamarsi} (\textquotedblleft{}to call oneself\textquotedblright{}), from Latin \textbf{clāmāre}, \emph{\textquotedblleft{}to cry out, to call.\textquotedblright{}} That root shouts through English:

\begin{itemize}
\raggedright
  \item \textbf{claim}, \textbf{exclaim}, \textbf{proclaim}, \textbf{acclaim}, \textbf{clamor} — all \textquotedblleft{}calling out.\textquotedblright{}
  \item Its Romance siblings all name people the same way: Spanish \textbf{llamarse} (\emph{me llamo}), Portuguese \textbf{chamar-se} (\emph{me chamo}) — the \emph{cl-} softening to \emph{ll-} (Spanish) or \emph{ch-} (Italian/Portuguese). One Latin \textquotedblleft{}call,\textquotedblright{} three naming verbs.
\end{itemize}
\end{cousinweb}

\begin{grammarlens}[title={Grammar Lens: the reflexive, and the dropping of \emph{io}}]
\emph{chiamo} alone is \textquotedblleft{}I call {[}someone else{]}\textquotedblright{}; \textbf{mi} chiamo is \textquotedblleft{}I call \textbf{myself}\textquotedblright{} — \textbf{reflexive}, the action looping back on you. The pronoun changes with the person: \textbf{mi} chiamo (I) $\to$ \textbf{ti} chiami (you) $\to$ \textbf{si} chiama (he/she/formal). And notice: no \emph{io} — the \emph{-o} ending already says \textquotedblleft{}I.\textquotedblright{} (You met that pro-drop last lesson.)
\end{grammarlens}

\subsection*{Your turn}
Say these aloud:

\begin{itemize}
\raggedright
  \item \textquotedblleft{}mi chiamo\textquotedblright{} — \emph{mee KYAH-moh}, hard \emph{k} in \emph{chiamo}
  \item \textquotedblleft{}Mi chiamo …\textquotedblright{} with your own name
  \item \textquotedblleft{}chiamo\textquotedblright{} then English \textquotedblleft{}claim, exclaim, clamor\textquotedblright{} — the \emph{clāmāre} family
\end{itemize}

\begin{tcolorbox}[breakable,colback=teal!4,colframe=teal!35!black,title={Before you move on}]
What does \emph{mi chiamo} literally mean? (\textquotedblleft{}I call myself.\textquotedblright{}) What Latin verb is \emph{chiamarsi} from, and its English cousins? (\emph{clāmāre} — claim, exclaim, clamor.) How is Italian \emph{ch} pronounced? (A hard \emph{k}.) Next: ask it back — \textbf{Come ti chiami?}
\end{tcolorbox}
\section[Come ti chiami?]{Come ti chiami? — \textquotedblleft{}what's your name?\textquotedblright{}}

\label{lesson:IT-C03-come-ti-chiami}



\begin{quote}
You have both pieces: \textbf{come} (\textquotedblleft{}how,\textquotedblright{} from your how-are-you chapter) and \textbf{chiamarsi} (\textquotedblleft{}to call oneself\textquotedblright{}). Like Spanish and French, Italian asks the name with \textbf{\textquotedblleft{}how\textquotedblright{}}, not \textquotedblleft{}what\textquotedblright{}: \emph{how do you call yourself?}
\end{quote}

\begin{cousinweb}
\begin{quote}
\textbf{Come ti chiami?} = \textquotedblleft{}What's your name?\textquotedblright{} (informal) literally: \textbf{\textquotedblleft{}How do you call yourself?\textquotedblright{}}
\end{quote}

\begin{itemize}
\raggedright
  \item \textbf{come} = \textquotedblleft{}how\textquotedblright{} ($\leftarrow$ \emph{quōmodo}; you met it in \emph{Come stai?}).
  \item \textbf{ti} = \textquotedblleft{}yourself\textquotedblright{} (the reflexive, matching \emph{mi} $\to$ \emph{ti}).
  \item \textbf{chiami} = \textquotedblleft{}you call\textquotedblright{} (the \emph{tu} form of \emph{chiamarsi}).
\end{itemize}

Two registers — the same \emph{tu / Lei} split you learned in Chapter 2:

\noindent
\begin{tabularx}{\linewidth}{@{}>{\raggedright\arraybackslash}X>{\raggedright\arraybackslash}X@{}}
\toprule
\textbf{register} & \textbf{question} \\
\midrule
\textbf{informal} (\emph{tu}) & \textbf{Come ti chiami?} \\
\textbf{formal} (\emph{Lei}) & \textbf{Come si chiama?} \\
\bottomrule
\end{tabularx}

Only the reflexive pronoun and the ending move: \emph{ti chiami} $\to$ \emph{si chiama} — the identical pattern to \emph{mi chiamo $\to$ si chiama}. The whole exchange:

\begin{quote}
— Come ti chiami? — Mi chiamo Marco. E tu?
\end{quote}

Note Italian asks the name with \emph{come} (\textquotedblleft{}how\textquotedblright{}), exactly like Spanish \emph{¿cómo se llama?} and French \emph{comment vous appelez-vous?} — the name is a matter of \emph{how} you're called, not \emph{what} you are.
\end{cousinweb}

\subsection*{Your turn}
Say these aloud:

\begin{itemize}
\raggedright
  \item \textquotedblleft{}Come ti chiami?\textquotedblright{} — informal, for a peer
  \item \textquotedblleft{}Come si chiama?\textquotedblright{} — formal, for a stranger (Lei)
  \item the whole swap — \textquotedblleft{}Come ti chiami? — Mi chiamo … E tu?\textquotedblright{}
\end{itemize}

\begin{tcolorbox}[breakable,colback=teal!4,colframe=teal!35!black,title={Before you move on}]
What does \emph{Come ti chiami?} literally ask? (\textquotedblleft{}How do you call yourself?\textquotedblright{}) What changes for the formal \emph{Lei} version? (\emph{ti chiami} $\to$ \emph{si chiama}.) Which \textquotedblleft{}how\textquotedblright{} word is it built on? (\emph{come} $\leftarrow$ \emph{quōmodo}.) Next: the reply when you meet — \textbf{piacere}.
\end{tcolorbox}
\section[piacere]{piacere — \textquotedblleft{}pleased to meet you,\textquotedblright{} i.e. \textquotedblleft{}a pleasure\textquotedblright{}}

\label{lesson:IT-C03-piacere}



\begin{quote}
Names exchanged, you close the introduction with one warm word: \textbf{piacere} — literally \emph{\textquotedblleft{}a pleasure.\textquotedblright{}} It's the twin of a word you'd say in Portuguese, and it's pure Latin \emph{please}.
\end{quote}

\begin{sounds}
\begin{itemize}
\raggedright
  \item \texttt{soft-c} — \textbf{piacere} = \emph{pya-\textbf{CHEH}-reh}: here \emph{c} is before \emph{e}, so it's soft (\emph{ch}), and the \emph{ia} glides: \emph{pya-CHE-re}.
\end{itemize}
\end{sounds}

\begin{cousinweb}
\textbf{piacere} = \textquotedblleft{}a pleasure\textquotedblright{} (and, as a verb, \textquotedblleft{}to please\textquotedblright{}), from Latin \textbf{placēre}, \emph{\textquotedblleft{}to please, to be pleasing.\textquotedblright{}} Said on meeting someone, it's short for \emph{\textquotedblleft{}{[}it's{]} a pleasure {[}to meet you{]}.\textquotedblright{}}

That \textbf{placēre} root is thoroughly at home in English:

\begin{itemize}
\raggedright
  \item \textbf{please}, \textbf{pleasure}, \textbf{pleasant} — the direct line.
  \item \textbf{complacent} (\textquotedblleft{}pleased \emph{with} oneself\textquotedblright{}), \textbf{complaisant}, \textbf{placid} (pleased $\to$ calm).
\end{itemize}

And across the family, \textquotedblleft{}pleased to meet you\textquotedblright{} is built from the \emph{same} idea nearly everywhere:

\begin{itemize}
\raggedright
  \item Portuguese \textbf{prazer} ($\leftarrow$ \emph{placēre} too — its twin).
  \item French \textbf{enchanté} (a different image — \textquotedblleft{}enchanted\textquotedblright{}).
  \item Spanish \textbf{mucho gusto} (\textquotedblleft{}much pleasure,\textquotedblright{} but from \emph{gustus}, \textquotedblleft{}taste\textquotedblright{} — a different Latin root).
\end{itemize}
\end{cousinweb}

\begin{grammarlens}[title={Grammar Lens: a one-word courtesy}]
\emph{Piacere} stands alone — you just say it, with a small nod, as you shake hands. For a touch more you can say \emph{molto piacere} (\textquotedblleft{}much pleasure\textquotedblright{}) or \emph{piacere di conoscerti} (\textquotedblleft{}pleasure to know you\textquotedblright{}).
\end{grammarlens}

\subsection*{Your turn}
Say these aloud:

\begin{itemize}
\raggedright
  \item \textquotedblleft{}piacere\textquotedblright{} — \emph{pya-CHEH-reh}
  \item \textquotedblleft{}piacere\textquotedblright{} then English \textquotedblleft{}please, pleasure, placid\textquotedblright{} — the \emph{placēre} root
  \item Italian \emph{piacere} and Portuguese \emph{prazer} — twins from \emph{placēre}
\end{itemize}

\begin{tcolorbox}[breakable,colback=teal!4,colframe=teal!35!black,title={Before you move on}]
What does \emph{piacere} literally mean, and from what? (\textquotedblleft{}A pleasure,\textquotedblright{} $\leftarrow$ \emph{placēre} \textquotedblleft{}to please.\textquotedblright{}) Its English cousins? (Please, pleasure, placid, complacent.) Which Portuguese word is its twin? (\emph{Prazer}.) Next: put the whole introduction together.
\end{tcolorbox}
\section[Practice]{Practice — introducing yourself in Italian}

\label{lesson:IT-C03-practice}



\begin{quote}
No new words. \emph{io, mi chiamo, come ti chiami, piacere} now assemble into a real first meeting — run it formally and informally until the register swap is automatic. This slots between your greetings (Ch. 1–2) and farewells (Ch. 4).
\end{quote}

\subsection*{The exchange}
\textbf{Informal} (a peer):

\begin{quote}
— Ciao! \textbf{Come ti chiami?} — \textbf{Mi chiamo} Marco. E tu? — Giulia. \textbf{Piacere!} — Piacere!
\end{quote}

\textbf{Formal} (a stranger — \emph{Lei}):

\begin{quote}
— Buongiorno. \textbf{Come si chiama?} — Mi chiamo Marco Rossi. E Lei? — Rossi. \textbf{Molto piacere.}
\end{quote}

\begin{grammarlens}[title={What you've built this chapter}]
\begin{itemize}
\raggedright
  \item \textbf{io} — \textquotedblleft{}I\textquotedblright{} ($\leftarrow$ \emph{ego}), usually \textbf{dropped} — the verb ending carries it.
  \item \textbf{mi chiamo} — \textquotedblleft{}I call myself\textquotedblright{} (\emph{chiamarsi} $\leftarrow$ \emph{clāmāre}, \textquotedblleft{}call out\textquotedblright{}; kin of Spanish \emph{me llamo}, Portuguese \emph{me chamo}).
  \item \textbf{Come ti chiami? / Come si chiama?} — \textquotedblleft{}how do you call yourself?\textquotedblright{}, informal / formal.
  \item \textbf{piacere} — \textquotedblleft{}a pleasure\textquotedblright{} ($\leftarrow$ \emph{placēre}; twin of Portuguese \emph{prazer}).
\end{itemize}
\end{grammarlens}

\subsection*{Your turn}
Say these aloud:

\begin{itemize}
\raggedright
  \item the full informal meeting, both voices
  \item the same formally — \emph{ti chiami $\to$ si chiama}, \emph{tu $\to$ Lei}
  \item chain it all — greet, name, ask, how-are-you, goodbye: \textquotedblleft{}Ciao! Mi chiamo … Come ti chiami? … Come stai? … A domani!\textquotedblright{}
\end{itemize}

\emph{Twice through:} Run a whole first meeting start to finish.

\begin{tcolorbox}[breakable,colback=teal!4,colframe=teal!35!black,title={Before you move on}]
Why does Italian usually drop \emph{io}? (The verb ending says who.) What does \emph{mi chiamo} literally mean, and its root? (\textquotedblleft{}I call myself\textquotedblright{}; \emph{clāmāre}.) What does \emph{piacere} mean? (\textquotedblleft{}A pleasure.\textquotedblright{}) Italian now runs greet $\to$ introduce $\to$ how-are-you $\to$ goodbye, end to end.
\end{tcolorbox}
